\begin{figure*}[t]
    \centering
    \resizebox{\textwidth}{!}{%
        \begin{tikzpicture}[
    font=\sffamily\small,
    flow/.style={
        rectangle,
        rounded corners=2pt,
        draw=black,
        line width=0.7pt,
        align=center,
        minimum height=1.1cm,
        text width=2.8cm,
        fill=blue!5
    },
    decision/.style={
        diamond,
        aspect=2.0,
        draw=black,
        line width=0.7pt,
        align=center,
        inner sep=1.5pt,
        text width=2.3cm,
        fill=orange!10
    },
    terminal/.style={
        ellipse,
        draw=black,
        line width=0.7pt,
        align=center,
        minimum height=0.9cm,
        text width=1.8cm,
        fill=green!10
    },
    block/.style={
        ellipse,
        draw=black,
        line width=0.7pt,
        align=center,
        minimum height=0.9cm,
        text width=1.8cm,
        fill=red!10
    },
    arrow/.style={-{Latex[length=2.4mm]}, line width=0.7pt},
    label/.style={font=\sffamily\scriptsize, fill=white, inner sep=1pt}
]

\node[flow] (candidate) {Strategy-selected\\candidate Pod};
\node[decision, right=0.8cm of candidate] (scope) {Inside configured scope?};
\node[decision, right=0.95cm of scope] (metrics) {Metrics available and complete?};
\node[flow, right=0.9cm of metrics] (aggregate) {Aggregate per-container usage and effective requests};
\node[flow, below=1.0cm of aggregate] (evaluate) {Evaluate CPU and memory independently, applying missing-request policy};
\node[decision, left=1.0cm of evaluate] (busy) {Either dimension busy?};
\node[terminal, left=1.0cm of busy] (allow) {Allow eviction};
\node[block, below=0.8cm of busy] (deny) {Block eviction};
\node[terminal, above=1.0cm of scope] (outscope) {Allow eviction};
\node[block, above=1.0cm of metrics] (nometrics) {Block eviction};

\draw[arrow] (candidate) -- (scope);
\draw[arrow] (scope) -- node[label, above] {yes} (metrics);
\draw[arrow] (metrics) -- node[label, above] {yes} (aggregate);
\draw[arrow] (aggregate) -- (evaluate);
\draw[arrow] (evaluate) -- (busy);
\draw[arrow] (busy) -- node[label, above] {no} (allow);
\draw[arrow] (busy) -- node[label, right] {yes} (deny);
\draw[arrow] (scope) -- node[label, right] {no} (outscope);
\draw[arrow] (metrics) -- node[label, right] {no} (nometrics);

\end{tikzpicture}

    }
    \caption{Actual Usage Evictor pre-eviction decision flow. Request dimensions
    are evaluated independently, while unavailable or incomplete runtime
    metrics block the eviction.}
    \label{fig:actualusage_workflow}
\end{figure*}

\section{Actual Usage Evictor Implementation}
\label{sec:implementation}

\subsection{Plugin Position and Scope}

The Actual Usage Evictor implements the Descheduler \textit{EvictorPlugin}
interface. Its general \textit{Filter} method is a no-op that returns true;
runtime usage is evaluated only in \textit{PreEvictionFilter}. Candidate
selection therefore remains the responsibility of the active strategy. The
plugin observes the selected Pod immediately before eviction and either allows
the pipeline to continue or rejects that candidate.

\begin{table*}[t]
    \centering
    \caption{Actual Usage Evictor Parameters}
    \label{tab:actualusage_parameters}
    \small
    \renewcommand{\arraystretch}{1.2}
    \begin{tabular}{@{}p{3.9cm}p{1.5cm}p{1.6cm}p{7.9cm}@{}}
        \toprule
        \textbf{Parameter} & \textbf{Type} & \textbf{Default} & \textbf{Behavior} \\
        \midrule
        \textit{cpuUsageThreshold} & float & \(0.80\) &
        Blocks when actual CPU divided by effective requested CPU reaches the threshold. \\
        \textit{memoryUsageThreshold} & float & \(0.80\) &
        Blocks when actual memory divided by effective requested memory reaches the threshold. \\
        \textit{missingRequestPolicy} & string & Allow &
        Skips a missing request dimension under \textit{Allow}; rejects it under \textit{Block}. \\
        \textit{namespaces} & object & optional &
        Restricts evaluation with namespace include or exclude rules. \\
        \textit{labelSelector} & object & optional &
        Restricts evaluation using Kubernetes label matching. \\
        \bottomrule
    \end{tabular}
\end{table*}

Operators can restrict protection with namespace include or exclude rules and
a Kubernetes label selector. A Pod outside the configured scope passes without
a metrics query. When both controls are present, both must match for the
runtime check to run. This permits targeted protection of user-facing or
latency-sensitive workloads while leaving batch workloads unaffected.

Figure~\ref{fig:actualusage_workflow} shows the decision flow. For an in-scope
Pod, the plugin reads Pod metrics and matches each regular container in the
metric response with the Pod specification. Usage and effective requests are
aggregated only after every container has a complete CPU and memory
observation.

\subsection{Request-Relative Usage Model}

For Pod \(p\), let \(A_{p,c}\) and \(A_{p,m}\) denote aggregate actual CPU and
memory usage. Let \(R_{p,c}\) and \(R_{p,m}\) denote effective CPU and memory
requests aggregated across all regular containers. For a dimension in which
every container has a request, the ratios are

\[
\rho_{p,c}=\frac{A_{p,c}}{R_{p,c}},
\qquad
\rho_{p,m}=\frac{A_{p,m}}{R_{p,m}}.
\]

CPU is aggregated in millicores and memory in bytes. Effective requests are
read from the admitted Pod object, so defaults inserted by a
\textit{LimitRange} are included without a separate lookup.

Given CPU threshold \(\theta_c\) and memory threshold \(\theta_m\), the filter
blocks the eviction when either dimension reaches its threshold:

\[
\mathrm{Block}(p)
\iff
\rho_{p,c}\geq\theta_c
\;\lor\;
\rho_{p,m}\geq\theta_m.
\]

Consequently, both ratios must be below their thresholds for the candidate to
pass. The default \(\theta_c=\theta_m=0.80\) classifies activity relative to
the request baseline; it does not mean that the container is at \(80\%\) of
its runtime limit. A Pod with a \(50\,m\) CPU request and \(1000\,m\) limit can
produce a ratio far above one while legitimately bursting below its limit.
Limits are deliberately excluded because proximity to throttling or an
out-of-memory boundary is a different decision from activity relative to the
scheduler baseline.

The implementation applies this decision to both CPU and memory. The
experiments in Section~\ref{sec:evaluation}, however, generate foreground CPU
work and trigger the event from CPU samples; they validate the CPU decision
path only. The memory path was not evaluated independently.

\subsection{Missing Data and Safety Behavior}

Missing runtime evidence and missing request baselines have different
semantics. The plugin fails closed when the metrics client is unavailable,
Pod metrics cannot be read, or any regular container lacks a CPU or memory
observation. In these cases it cannot establish that eviction is safe.

If a container lacks an effective request for one resource, the entire
Pod-level dimension is marked missing rather than dividing by a partial
denominator. The dimension is then resolved using
\textit{missingRequestPolicy}. The default \textit{Allow} skips only that
dimension; the other dimension can still block through the OR condition.
\textit{Block} treats the missing dimension as busy. Thus, \textit{Allow} does
not claim that a request-less Pod is idle---it states that request-relative
busyness cannot be inferred for that dimension.

\subsection{Configuration}

Table~\ref{tab:actualusage_parameters} summarizes the public plugin arguments.
The complete decision procedure and an example policy are provided in
Appendices~\ref{appendix:actualusage_algorithm}
and~\ref{appendix:actualusage_config}.
